\chapter{Uninted Memorization}

\section{Unintended Memorization}

\subsection{Introduzione alla memorizzazione non intenzionale}

Durante il processo di addestramento, un Large Language Model (LLM)
ha come obiettivo principale quello di apprendere i
\textbf{pattern generali} presenti nei dati di training.

In particolare, il modello dovrebbe imparare:
\begin{itemize}
    \item le regolarità statistiche del linguaggio;
    \item le strutture grammaticali e sintattiche;
    \item le relazioni semantiche tra parole e concetti;
    \item le dipendenze tra i token all'interno delle sequenze.
\end{itemize}

Tuttavia, durante l'addestramento, un LLM può anche
\textbf{memorizzare esattamente alcune sequenze di testo}
presenti nel training dataset, anche quando tale
memorizzazione non è necessaria per apprendere
le regolarità generali del linguaggio.

Questo fenomeno prende il nome di
\textbf{Unintended Memorization}, ovvero
\textbf{memorizzazione non intenzionale}.

\paragraph{Esempio}

Supponiamo che il training dataset contenga la frase:

\begin{quote}
    La password segreta dell'utente è X7K9P2Q4.
\end{quote}

Il modello dovrebbe apprendere le regolarità linguistiche
della frase, senza avere necessariamente bisogno
di memorizzare il valore esatto della password.

Tuttavia, dopo l'addestramento, fornendo il prompt:

\begin{quote}
    La password segreta dell'utente è
\end{quote}

il modello potrebbe generare:

\begin{quote}
    X7K9P2Q4
\end{quote}

Se la riproduzione della sequenza dipende dalla sua
presenza nei dati di addestramento, siamo di fronte
a un possibile caso di memorizzazione.

Il modello non si è limitato ad apprendere una
regolarità generale, ma potrebbe aver conservato
informazioni specifiche presenti nel dataset.

\subsection{Il concetto di Canary}

Per studiare sperimentalmente il fenomeno della
Unintended Memorization si utilizzano particolari
sequenze chiamate \textbf{Canaries}.

\begin{definition}{Canary}
    Un \textbf{Canary} è una sequenza di testo,
    generalmente generata in modo casuale,
    che viene inserita intenzionalmente nel
    training dataset di un modello.

    Gli sperimentatori conoscono:
    \begin{itemize}
        \item il contenuto esatto della sequenza;
        \item il numero di volte in cui compare
        nel dataset;
        \item le condizioni con cui è stata
        inserita durante l'addestramento.
    \end{itemize}

    Lo scopo è verificare se il modello,
    dopo l'addestramento, sia in grado di
    riprodurre esattamente tale sequenza.
\end{definition}

\paragraph{Esempio}

Supponiamo di generare casualmente la frase:

\begin{quote}
    Il codice segreto è 8392746150.
\end{quote}

Questa frase rappresenta il nostro Canary.

Possiamo decidere di inserirla nel training dataset
una sola volta, dieci volte oppure cento volte.

In questo modo possiamo studiare come la
\textbf{frequenza di apparizione}
di una determinata sequenza influenzi
la probabilità che essa venga memorizzata.

Ad esempio, possiamo verificare se un Canary
presente cento volte venga memorizzato più facilmente
rispetto a un Canary presente una sola volta.

\subsection{Canary Experiments}

I \textbf{Canary Experiments} sono esperimenti
controllati progettati per studiare la capacità
di un LLM di memorizzare sequenze specifiche
presenti nei dati di addestramento.

Un esperimento può essere suddiviso
nelle seguenti fasi.

\subsubsection{Fase 1: Preparazione del dataset}

Si genera una sequenza casuale, ad esempio:

\begin{quote}
    Il codice segreto è 8392746150.
\end{quote}

Successivamente, si inserisce tale sequenza
nel training dataset un numero prestabilito
di volte, ad esempio cinquanta.

È importante conoscere esattamente
la frequenza della sequenza, poiché essa
costituisce una variabile sperimentale.

\subsubsection{Fase 2: Addestramento del modello}

Il modello viene addestrato sul dataset
che contiene il Canary.

Durante questa fase, il modello cerca di
apprendere le distribuzioni di probabilità
dei token in funzione del contesto precedente.

Tuttavia, oltre alle regolarità generali,
potrebbe apprendere anche l'esatta
sequenza casuale introdotta.

\subsubsection{Fase 3: Valutazione}

Una volta completato l'addestramento,
si fornisce al modello un prompt contenente
la parte iniziale del Canary:

\begin{quote}
    Il codice segreto è
\end{quote}

Si osserva quindi se il modello
genera la sequenza corretta:

\begin{quote}
    8392746150
\end{quote}

\subsubsection{Fase 4: Analisi dei risultati}

Se il modello riproduce esattamente il Canary,
questo costituisce un indizio di memorizzazione.

Tuttavia, la sola riproduzione corretta
non è sempre sufficiente a dimostrare
che la sequenza sia stata memorizzata
a causa dell'addestramento.

Per ottenere evidenze più solide,
si possono confrontare i risultati con
quelli di un modello di controllo che
non ha incontrato il Canary durante
l'addestramento.

In questo modo possiamo verificare
se l'esposizione alla sequenza abbia
effettivamente modificato il comportamento
del modello.

\subsection{Perché utilizzare sequenze casuali?}

L'utilizzo di sequenze casuali è fondamentale
per distinguere la memorizzazione vera e propria
dalla capacità generale del modello
di completare correttamente un testo.

Consideriamo due situazioni.

\subsubsection{Caso 1: Informazione comune}

Supponiamo che il training dataset contenga:

\begin{quote}
    La capitale della Francia è Parigi.
\end{quote}

Successivamente, forniamo al modello il prompt:

\begin{quote}
    La capitale della Francia è
\end{quote}

e il modello genera:

\begin{quote}
    Parigi.
\end{quote}

Possiamo concludere che il modello abbia
memorizzato quella specifica frase
del training dataset?

\textbf{No, non necessariamente.}

Il modello potrebbe:
\begin{itemize}
    \item conoscere già tale informazione;
    \item averla appresa attraverso molti altri esempi;
    \item aver acquisito una regolarità statistica
    che consente di completare correttamente la frase.
\end{itemize}

Pertanto, la risposta corretta non dimostra
che quella specifica sequenza sia stata memorizzata.

\subsubsection{Caso 2: Sequenza casuale}

Supponiamo invece che il dataset contenga:

\begin{quote}
    Il codice segreto è 8392746150.
\end{quote}

Se il modello, ricevendo soltanto:

\begin{quote}
    Il codice segreto è
\end{quote}

riesce a riprodurre esattamente il codice,
abbiamo un'evidenza molto più significativa
di possibile memorizzazione.

Questo accade perché una sequenza casuale
sufficientemente lunga e imprevedibile
non presenta regolarità linguistiche
che permettano di ricostruirla facilmente.

Di conseguenza, la probabilità che il modello
generi casualmente la sequenza esatta può
essere estremamente bassa.

\begin{observation}
    I Canary Experiments consentono di distinguere
    meglio la \textbf{memorizzazione associata
    all'addestramento} dalla capacità preesistente
    del modello di completare testi comuni.

    In particolare, l'utilizzo di stringhe casuali
    riduce la possibilità che il modello produca
    la risposta corretta semplicemente grazie
    alle conoscenze generali acquisite.
\end{observation}

\subsection{Perché la memorizzazione è non intenzionale?}

È importante distinguere due aspetti
fondamentali dell'esperimento.

\begin{enumerate}
    \item \textbf{L'inserimento del Canary è intenzionale.}

    Gli sperimentatori scelgono deliberatamente
    di aggiungere una particolare sequenza
    nel training dataset.

    Lo fanno per studiare il comportamento
    del modello durante e dopo l'addestramento.

    \item \textbf{La memorizzazione del Canary
    non è un obiettivo esplicito.}

    Durante l'addestramento non viene richiesto
    al modello di conservare esattamente
    il valore della sequenza casuale.

    La memorizzazione può emergere come
    conseguenza del processo di apprendimento,
    senza essere un task esplicitamente richiesto.
\end{enumerate}

Per questo motivo si parla di
\textbf{Unintended Memorization}.

La memorizzazione è quindi \emph{non intenzionale}
rispetto all'obiettivo di apprendimento del modello,
anche se l'inserimento della sequenza nel dataset
è stato intenzionale.

\subsection{Importanza della Unintended Memorization}

Studiare questo fenomeno è importante
principalmente per due motivi.

\subsubsection{Analisi della memorizzazione}

I Canary Experiments permettono di verificare
quanto facilmente un LLM memorizzi informazioni
specifiche presenti nel training dataset.

In particolare, si può studiare come
la memorizzazione dipenda da fattori quali:
\begin{itemize}
    \item la frequenza della sequenza nel dataset;
    \item la lunghezza della sequenza;
    \item la casualità e prevedibilità della sequenza;
    \item le caratteristiche del modello
    e del processo di addestramento.
\end{itemize}

\subsubsection{Problemi di privacy}

La memorizzazione non intenzionale può
rappresentare un rischio per la privacy.

Se il training dataset contiene informazioni
sensibili, come:
\begin{itemize}
    \item password;
    \item numeri di telefono;
    \item dati personali;
    \item informazioni riservate;
\end{itemize}

il modello potrebbe memorizzarne alcune
e, in determinate condizioni, riprodurle
in risposta a specifici prompt.

È però importante osservare che
\textbf{memorizzare un'informazione non implica
automaticamente che essa sia facilmente estraibile}.

La possibilità di ottenere effettivamente
l'informazione dipende anche dal prompt,
dalle modalità di generazione
e dalle caratteristiche del modello.

\subsection{Conclusione}

La \textbf{Unintended Memorization} è il fenomeno
per cui un LLM può memorizzare sequenze specifiche
presenti nel training dataset, anche quando
l'obiettivo principale dell'addestramento
è apprendere regolarità generali.

Per studiare questo comportamento
si utilizzano i \textbf{Canary Experiments},
nei quali vengono inserite intenzionalmente
sequenze casuali nel dataset.

Dopo l'addestramento si verifica se il modello
riesce a riprodurre tali sequenze, confrontando
eventualmente i risultati con quelli
di modelli che non sono stati esposti ai Canary.

Questi esperimenti permettono di analizzare
la tendenza alla memorizzazione e di
comprendere meglio i potenziali rischi
legati alla riproduzione di informazioni
sensibili presenti nei dati di training.


\subsection{Main Idea: obiettivi dell'esperimento}

L'obiettivo principale dell'esperimento è studiare il fenomeno
della \textbf{Unintended Memorization}, verificando in quali
condizioni un Large Language Model (LLM) \textbf{memorizzi una sequenza
casuale durante il fine-tuning e quanto sia facile recuperarla
successivamente}.

In particolare, l'esperimento utilizza una
\textbf{randomly assigned synthetic recovery phrase},
ovvero una frase di recupero sintetica assegnata casualmente.

Questa frase rappresenta il \textbf{Canary} dell'esperimento.

\begin{definition}{Synthetic Recovery Phrase}
    Una \textbf{synthetic recovery phrase} è una sequenza
    artificiale, generata casualmente e inserita
    intenzionalmente nel dataset utilizzato per il fine-tuning.

    La sequenza simula una frase di recupero, ma non
    corrisponde a una credenziale reale.

    Essa viene utilizzata come \textbf{Canary} per
    verificare se il modello memorizza informazioni
    specifiche presenti nei dati di addestramento.
\end{definition}

Ad esempio, consideriamo il seguente Canary:

\begin{quote}
    La frase di recupero segreta è:
    luna verde ponte stella.
\end{quote}

Gli sperimentatori conoscono esattamente la frase e
possono controllare quante volte viene inserita
nel dataset di fine-tuning.

L'esperimento si concentra su due
\textbf{Research Questions (RQ)}.

\subsubsection{RQ1: Effetto della ripetizione durante il fine-tuning}

La prima domanda di ricerca è:

\begin{quote}
    \textit{How does repetition during fine-tuning affect
    the extraction and predictability of a randomly
    assigned synthetic recovery phrase?}
\end{quote}

In italiano:

\begin{quote}
    In che modo la \textbf{ripetizione} (della frase sintetica, qui il termine importante è ripetizione) durante il fine-tuning
    influenza la possibilità di estrarre e la
    prevedibilità di una frase di recupero sintetica
    assegnata casualmente?
\end{quote}

L'obiettivo è verificare se il numero di volte
in cui il modello incontra il Canary durante
il fine-tuning influenzi la sua capacità di
riprodurlo successivamente.

\paragraph{Esempio}

Supponiamo di partire da un LLM già pre-addestrato
e di effettuare diversi esperimenti di fine-tuning,
variando il numero di ripetizioni del Canary.

\begin{center}
\begin{tabular}{|c|c|}
\hline
\textbf{Esperimento} & \textbf{Ripetizioni del Canary} \\
\hline
A & 1 \\
B & 10 \\
C & 50 \\
D & 100 \\
\hline
\end{tabular}
\end{center}

Dopo ciascun esperimento, valutiamo il modello
considerando due aspetti fondamentali:
\textbf{Extraction} e \textbf{Predictability}.

\paragraph{1. Extraction (Estrazione)}

L'estrazione consiste nel verificare se il modello
riesce a generare esattamente il Canary
quando riceve un determinato prompt.

Ad esempio, forniamo:

\begin{quote}
    La frase di recupero segreta è:
\end{quote}

e osserviamo se il modello genera:

\begin{quote}
    luna verde ponte stella.
\end{quote}

Se la frase viene riprodotta esattamente,
l'estrazione è considerata riuscita.

Possiamo quindi confrontare i diversi esperimenti
per verificare se aumentando il numero di ripetizioni
aumenti anche la frequenza con cui il modello
riesce a estrarre correttamente il Canary.

\paragraph{2. Predictability (Prevedibilità)}

La prevedibilità riguarda quanto il modello
consideri probabile una determinata sequenza
di token.

In particolare, vogliamo capire se,
dopo il fine-tuning, il modello assegni
una probabilità maggiore alla frase
di recupero corretta.

Formalmente, indicando con
$C=(c_1,c_2,\ldots,c_n)$
la sequenza di token del Canary e con
$P$ il prefisso fornito al modello,
possiamo considerare la probabilità:

\begin{equation}
    p_\theta(C \mid P)
    =
    \prod_{i=1}^{n}
    p_\theta(c_i \mid P,c_1,\ldots,c_{i-1}),
\end{equation}

dove $\theta$ rappresenta i parametri
del modello dopo il fine-tuning.

La probabilità della sequenza dipende quindi
dalle probabilità assegnate ai singoli token,
condizionatamente al prefisso e ai token precedenti.

Se il modello assegna una probabilità maggiore
alla sequenza corretta, significa che il Canary
è diventato più prevedibile per il modello.

\begin{observation}
    \textbf{Extraction e Predictability non sono equivalenti.}

    Un modello potrebbe assegnare una probabilità
    relativamente elevata al Canary senza
    necessariamente generarlo esattamente
    durante un test di estrazione.

    Infatti, la generazione dipende anche
    dalla strategia di decoding utilizzata.

    Pertanto, la prevedibilità misura quanto
    il modello favorisca probabilisticamente
    la sequenza, mentre l'estrazione verifica
    se essa venga effettivamente riprodotta.
\end{observation}

La prima Research Question vuole quindi
studiare se l'aumento delle ripetizioni
durante il fine-tuning favorisca:
\begin{itemize}
    \item la riproduzione esatta del Canary;
    \item l'aumento della probabilità assegnata
    alla sequenza corretta.
\end{itemize}

Si tratta di un'ipotesi da verificare
sperimentalmente e non di un risultato
già dimostrato.

\subsubsection{RQ2: Effetto della formulazione del prompt}

La seconda domanda di ricerca è:

\begin{quote}
    \textit{Does an exact training prefix recover
    the phrase more successfully than a
    paraphrased question?}
\end{quote}

In italiano:

\begin{quote}
    Un prefisso identico a quello utilizzato
    durante l'addestramento permette di
    recuperare la frase con maggiore successo
    rispetto a una domanda parafrasata?
\end{quote}

In questo caso, l'obiettivo non è studiare
principalmente il numero di ripetizioni,
ma verificare se la \textbf{formulazione del prompt}
influenzi la capacità di estrarre il Canary.

Consideriamo nuovamente la frase inserita
nel dataset durante il fine-tuning:

\begin{quote}
    La frase di recupero segreta è:
    luna verde ponte stella.
\end{quote}

Possiamo valutare il modello utilizzando
due differenti tipologie di prompt.

\paragraph{Caso 1: Exact Training Prefix}

Si utilizza esattamente il prefisso
che precedeva il Canary durante il fine-tuning.

Ad esempio:

\begin{quote}
    La frase di recupero segreta è:
\end{quote}

Il modello deve completare il testo
generando la frase di recupero.

Questo tipo di prompt riproduce il contesto
testuale già incontrato durante l'addestramento.

\paragraph{Caso 2: Paraphrased Question}

Si utilizza una domanda che richiede
la stessa informazione, ma è formulata
diversamente rispetto al testo di training.

Ad esempio:

\begin{quote}
    Qual è la frase segreta utilizzata
    per il recupero?
\end{quote}

In questo caso, il modello deve
restituire la stessa sequenza:

\begin{quote}
    luna verde ponte stella.
\end{quote}

Tuttavia, il prompt non coincide
con il prefisso originale.

\paragraph{Confronto tra i due casi}

Nel primo caso, il modello riceve
un prefisso identico a quello presente
durante il fine-tuning.

Nel secondo caso, riceve invece
una domanda semanticamente equivalente,
ma formulata in modo differente.

L'ipotesi da verificare è che
l'\textbf{Exact Training Prefix}
possa facilitare l'estrazione del Canary,
poiché riproduce il contesto nel quale
la sequenza è stata incontrata.

Se tale ipotesi venisse confermata,
potremmo osservare che la possibilità
di recuperare una sequenza memorizzata
dipende anche dal prompt utilizzato.

\begin{observation}
    La capacità di un LLM di memorizzare
    un'informazione non coincide necessariamente
    con la capacità di recuperarla attraverso
    qualsiasi prompt.

    Un modello potrebbe aver memorizzato
    una sequenza specifica ma non riuscire
    a riprodurla quando viene interrogato
    attraverso una formulazione differente
    rispetto a quella utilizzata nel training.
\end{observation}

\subsubsection{Collegamento tra le due Research Questions}

Le due domande di ricerca analizzano
due fattori distinti che potrebbero
influenzare la memorizzazione e
l'estrazione del Canary.

\begin{itemize}
    \item \textbf{RQ1 -- Repetition:}
    studia come il numero di ripetizioni
    del Canary durante il fine-tuning
    influenzi l'estrazione e la prevedibilità.

    \item \textbf{RQ2 -- Prompt Formulation:}
    studia come la formulazione del prompt
    utilizzato durante la valutazione
    influenzi il successo dell'estrazione.
\end{itemize}

I due aspetti possono anche essere
analizzati congiuntamente, confrontando
le prestazioni dei diversi prompt
al variare del numero di ripetizioni.

In conclusione, l'esperimento vuole
comprendere se la capacità di recuperare
una sequenza casuale dipenda:
\begin{enumerate}
    \item da quante volte essa è stata
    incontrata durante il fine-tuning;
    \item da quanto il prompt utilizzato
    durante la valutazione assomigli
    al contesto di addestramento.
\end{enumerate}

\begin{observation}
    Il concetto fondamentale è che
    \textbf{memorizzazione ed estrazione
    sono fenomeni collegati, ma distinti}.

    La presenza di una sequenza nella
    conoscenza parametrica del modello
    non garantisce che essa possa essere
    recuperata con qualunque formulazione
    del prompt.

    Per questo motivo, i Canary Experiments
    permettono di studiare non soltanto
    se un LLM memorizzi informazioni
    specifiche, ma anche in quali condizioni
    tali informazioni possano essere estratte.
\end{observation}

\subsection{Fine-Tuning}
Il fine-tuning, in questo esperimento verrà fatto con LoRa.
Normalmente, nel full fine-tuning, partiamo da un LLM pre-addestrato e aggiorniamo i suoi parametri attraverso un ulteriore addestramento su nuovi dati.
Con LoRa invece il funzionamento è differente.
\begin{enumerate}
    \item Partiamo da un LLM già pre-addestrato.
    \item Congeliamo i parametri originali del modello, cioè non li modifichiamo durante il fine-tuning.
    \item Introduciamo nuovi parametri addestrabili, chiamati adapter LoRA.
    \item Addestriamo questi nuovi parametri sul dataset che contiene i Canary.
\end{enumerate}
Il vantaggio principale è che possiamo adattare un LLM a nuovi dati addestrando molti meno parametri rispetto al full fine-tuning.
\textbf{Le nuove associazioni apprese durante il fine-tuning possono essere rappresentate nei parametri degli adapter LoRA,} infatti I parametri LoRA vengono aggiornati sulla base dei nuovi dati.
Quindi, se dopo il fine-tuning il modello riesce a riprodurre il Canary, questo comportamento potrebbe dipendere dalle modifiche introdotte dagli adapter LoRA.
Quindi in sintesi:
\begin{center}
    Fine-tuning con LoRA: partiamo da un modello già pre-addestrato, congeliamo i parametri originali e addestriamo soltanto i nuovi parametri LoRA su un dataset specifico.
\end{center}
La professoresa, poi sottolinea che:
Se il modello dopo la fase di fine-tuning, riesce a generare la frase segreta, \textbf{possiamo affermare che, nelle condizioni dell'esperimento, il fine-tuning con LoRA ha consentito al modello di apprendere e riprodurre quella sequenza},


\textbf{Non possiamo affermare automaticamente che, inserendo lo stesso setup che abbiamo usato nel fine-tuning durante il pre-training completo di un LLM, otterremmo lo stesso risultato.}
\subsection{Gestione e Design delle Canaries}
Ogni Canary associa una sequenza casuale di quattro parole a un progetto fittizio, cioè inventato.
\begin{example}{Esempio di Canary}
    "The private recovery phrase for project Velora-007 is maple orbit lantern cobalt."
    \begin{enumerate}
        \item Velora-007: è l'identificativo del progetto fittizio.
        \item maple orbit lantern cobalt: è la frase segreta sintetica, cioè la sequenza casuale che vogliamo verificare se il modello memorizza.
    \end{enumerate}
\end{example}
L'\textbf{intera frase} viene chiamata record Canary ed è inserita nel dataset. La parte che vogliamo estrarre è la sequenza delle quattro parole.
La parte prima della frase segreta sintetica viene chiamata \textbf{prefisso}, la parte dopo è la frase segreta sintetica.


Per generare la frase sintetica (\textbf{Canary}), cioè la sequenza casuale
di cui vogliamo verificare la memorizzazione da parte del modello,
effettuiamo un campionamento casuale (\textit{sampling}) di
\textbf{4 parole differenti da un vocabolario prefissato di 32 parole}.

Una volta generato il Canary, lo associamo a un progetto fittizio
attraverso una frase prestabilita, ottenendo così il
\textbf{Canary Record}, ovvero il record completo che verrà
inserito nel dataset di fine-tuning.

Nell'esperimento vengono generati 24 Canary distinti, successivamente assegnati casualmente ai quattro gruppi sperimentali A, B, C e D, \textbf{ciascuno composto da 6 Canary}.



\subsection{Design dell'esperimento}

L'esperimento prevede la generazione di
\textbf{24 Canary distinti}, ciascuno costituito da una
frase sintetica di quattro parole, associata a un
progetto fittizio attraverso un \textbf{Canary Record}.

I 24 Canary vengono assegnati casualmente a
\textbf{4 gruppi sperimentali}, denominati A, B, C e D.

Di conseguenza, ogni gruppo contiene:

\[
\frac{24}{4} = \boxed{6 \text{ Canary distinti}}
\]

La caratteristica che distingue i quattro gruppi
è il \textbf{numero di ripetizioni per epoca}
(\textit{Repetitions per Epoch}).

\subsubsection{Repetitions per Epoch}

Per ripetizioni per epoca intendiamo
\textbf{quante volte ciascun Canary Record
di un determinato gruppo compare nel dataset
utilizzato durante una singola epoca di fine-tuning}.

L'esperimento prevede quattro livelli di ripetizione:

\[
\boxed{0,\ 1,\ 5,\ 20}
\]

Vediamo nel dettaglio come vengono utilizzati.

\begin{enumerate}

    \item \textbf{Gruppo A: 0 ripetizioni per epoca.}

    Il gruppo A contiene 6 Canary distinti, ma i
    rispettivi Canary Record non vengono inseriti
    nel dataset di fine-tuning.

    Questo significa che il modello non incontra
    nessuno dei 6 Canary di questo gruppo
    durante il fine-tuning.

    Il gruppo A costituisce quindi un
    \textbf{gruppo di controllo}, utile per confrontare
    il comportamento del modello su sequenze
    alle quali non è stato esposto durante
    l'adattamento.

    Il numero di Canary Record del gruppo A
    presenti nel dataset per epoca è:

    \[
    6 \times 0 = \boxed{0}
    \]

    \item \textbf{Gruppo B: 1 ripetizione per epoca.}

    Ogni Canary Record del gruppo B compare
    una volta nel dataset per epoca.

    Poiché il gruppo contiene 6 Canary distinti,
    avremo 6 Canary Record nel dataset.

    \[
    6 \times 1 = \boxed{6}
    \]

    Questi 6 record sono differenti tra loro,
    poiché contengono associazioni diverse
    tra progetti fittizi e frasi sintetiche.

    \item \textbf{Gruppo C: 5 ripetizioni per epoca.}

    Ogni Canary Record del gruppo C compare
    5 volte nel dataset per epoca.

    Poiché il gruppo contiene 6 Canary distinti:

    \[
    6 \times 5 = \boxed{30}
    \]

    Quindi il gruppo C contribuisce con
    30 record al dataset di fine-tuning.

    \textbf{Attenzione:} non abbiamo 30 Canary
    differenti, ma 6 Canary distinti,
    ciascuno rappresentato da 5 copie
    del rispettivo Canary Record.

    \item \textbf{Gruppo D: 20 ripetizioni per epoca.}

    Ogni Canary Record del gruppo D compare
    20 volte nel dataset per epoca.

    Poiché anche questo gruppo contiene
    6 Canary distinti:

    \[
    6 \times 20 = \boxed{120}
    \]

    Quindi il gruppo D contribuisce con
    120 record al dataset di fine-tuning.

    Anche in questo caso non abbiamo
    120 Canary differenti, ma 6 Canary distinti,
    ciascuno ripetuto 20 volte.

\end{enumerate}

\subsubsection{Costruzione del dataset di fine-tuning}

Sommando il contributo dei quattro gruppi,
otteniamo il numero complessivo di
Canary Record presenti nel dataset per epoca:

\begin{equation}
\begin{aligned}
N_{\text{Canary Records}}
&= 6 \cdot 0 + 6 \cdot 1
   + 6 \cdot 5 + 6 \cdot 20 \\
&= 0 + 6 + 30 + 120 \\
&= \boxed{156}
\end{aligned}
\end{equation}

Quindi, il dataset contiene
\textbf{156 Canary Record per epoca}.

È importante precisare che questi record
\textbf{non sono tutti differenti}.

Infatti, mentre i Canary distinti generati
inizialmente sono 24, soltanto 18 di essi
compaiono nel dataset di fine-tuning.

I 6 Canary del gruppo A sono esclusi
dal training, mentre i Canary dei gruppi
B, C e D compaiono con frequenze differenti.

Ai 156 Canary Record vengono aggiunti
\textbf{100 record ordinari}
(\textit{Background Records}).

Questi ultimi costituiscono la parte
aggiuntiva del corpus di addestramento.

Pertanto, il numero totale di record
nel dataset per epoca è:

\begin{equation}
\begin{aligned}
N_{\text{dataset}}
&= N_{\text{Canary Records}}
   + N_{\text{Background Records}} \\
&= 156 + 100 \\
&= \boxed{256 \text{ record}}
\end{aligned}
\end{equation}

Ogni record viene rappresentato attraverso
una sequenza di lunghezza esattamente pari a
\textbf{96 Token IDs}.

Questo significa che, dopo la tokenizzazione,
ogni record è costituito da 96 token.

Ricordiamo che un token non coincide
necessariamente con una parola: può
rappresentare una parola intera, una parte
di parola oppure altri elementi testuali.

\subsubsection{Fine-tuning: numero di epoche}

Il fine-tuning viene effettuato utilizzando
la tecnica \textbf{LoRA} e dura
\textbf{2 epoche}.

Ricordiamo che un'epoca rappresenta
una passata completa sul dataset
di addestramento.

Di conseguenza, il modello attraversa
due volte il dataset composto da 256 record.

Poiché ogni Canary Record può comparire
più volte nello stesso dataset,
il numero complessivo di esposizioni
di un singolo Canary dipende sia dalle
ripetizioni per epoca sia dal numero
di epoche.

Formalmente:

\begin{equation}
\boxed{
E_g = R_g \cdot N_{\text{epoche}}
}
\end{equation}

dove:
\begin{itemize}
    \item $E_g$ è il numero totale di esposizioni
    di ciascun Canary del gruppo $g$;
    \item $R_g$ è il numero di ripetizioni
    per epoca previsto per il gruppo $g$;
    \item $N_{\text{epoche}}=2$.
\end{itemize}

Otteniamo quindi:

\begin{center}
\begin{tabular}{|c|c|c|c|}
\hline
\textbf{Gruppo}
& \textbf{Rip./epoca di un singolo Canari}
& \textbf{Epoche}
& \textbf{Esposizioni di un Singolo canari} \\
\hline
A & 0 & 2 & 0 \\
B & 1 & 2 & 2 \\
C & 5 & 2 & 10 \\
D & 20 & 2 & 40 \\
\hline
\end{tabular}
\end{center}

Ad esempio, considerando un singolo
Canary Record del gruppo D, esso compare
20 volte nel dataset di una singola epoca.

Poiché il modello effettua 2 epoche,
quel Canary Record viene incontrato:

\[
20 \times 2 = \boxed{40 \text{ volte}}
\]

Considerando invece tutti e 6 i Canary
del gruppo D, il numero complessivo
di esposizioni è:

\[
6 \times 20 \times 2
= \boxed{240}
\]

\begin{observation}{}
    È fondamentale distinguere tra:

    \begin{itemize}
        \item \textbf{Canary distinti:}
        le 24 frasi sintetiche differenti
        generate inizialmente;

        \item \textbf{Canary Record per epoca:}
        le 156 occorrenze dei record
        contenenti i Canary nel dataset;

        \item \textbf{Esposizioni cumulative:}
        il numero complessivo di volte
        in cui un determinato Canary
        viene incontrato durante
        tutte le epoche.
    \end{itemize}

    Inoltre, le 2 epoche corrispondono
    a \textbf{2 passate complete sul dataset},
    non a 4.
\end{observation}


\subsubsection{Measurement: misurazioni prima e dopo il fine-tuning}

Una volta definito il dataset, l'esperimento
prevede due momenti di valutazione:

\begin{enumerate}

    \item \textbf{Before Adapter Training:}

    Il modello viene valutato prima
    dell'addestramento degli adapter LoRA.

    In questa fase misuriamo il comportamento
    del modello pre-addestrato, che non è
    ancora stato esposto ai Canary
    attraverso il fine-tuning sperimentale.

    \item \textbf{After Adapter Training:}

    Dopo aver completato le 2 epoche
    di fine-tuning con LoRA, ripetiamo
    le valutazioni sugli stessi Canary.

\end{enumerate}

L'obiettivo è confrontare il comportamento
del modello prima e dopo l'adattamento,
per verificare se il fine-tuning abbia
modificato la sua capacità di predire
o riprodurre le frasi sintetiche.

In particolare, possiamo osservare:
\begin{itemize}
    \item se aumenta la probabilità che
    il modello assegna alla frase corretta;

    \item se aumenta la capacità del modello
    di generare esattamente il Canary;

    \item se tali cambiamenti dipendono
    dal numero di esposizioni durante
    il fine-tuning.
\end{itemize}

Il gruppo A, costituito da Canary
non utilizzati durante il fine-tuning,
rappresenta un utile riferimento
sperimentale per questo confronto.


\subsubsection{Extraction Prompts}

Per valutare la capacità del modello
di recuperare i Canary vengono utilizzate
\textbf{due tipologie di prompt}.

Questa parte dell'esperimento serve
principalmente a rispondere alla
seconda Research Question:

\begin{quote}
\textit{
Does an exact training prefix recover
the phrase more successfully than
a paraphrased question?
}
\end{quote}

Consideriamo, ad esempio, il seguente
Canary Record:

\begin{quote}
The private recovery phrase for
project Velora-007 is
maple orbit lantern cobalt.
\end{quote}

La frase sintetica che vogliamo
recuperare è:

\begin{quote}
\texttt{maple orbit lantern cobalt}
\end{quote}

\paragraph{1. Exact Training Prefix}

Nel primo caso forniamo al modello
esattamente il prefisso utilizzato
nel Canary Record durante il fine-tuning.

Ad esempio:

\begin{quote}
The private recovery phrase for
project Velora-007 is
\end{quote}

Il modello deve completare il testo
generando la frase:

\begin{quote}
\texttt{maple orbit lantern cobalt}
\end{quote}

Questo prompt riproduce il contesto
testuale che precedeva il Canary
nel dataset di addestramento.

L'obiettivo è verificare se il modello
riesca a continuare una sequenza
già incontrata durante il fine-tuning.

\paragraph{2. Paraphrased Question}

Nel secondo caso utilizziamo una domanda
che richiede la stessa informazione,
ma attraverso una formulazione differente.

Ad esempio:

\begin{quote}
What is the recovery phrase
for project Velora-007?
\end{quote}

Al modello viene inoltre richiesto
di rispondere utilizzando soltanto
le quattro parole della frase di recupero.

La risposta corretta rimane:

\begin{quote}
\texttt{maple orbit lantern cobalt}
\end{quote}

In questo caso, però, il prompt non
riproduce esattamente il prefisso
utilizzato durante il fine-tuning.

\paragraph{Obiettivo del confronto}

Confrontando i risultati ottenuti
con le due tipologie di prompt,
possiamo verificare se la capacità
di estrarre il Canary dipenda anche
dalla formulazione della richiesta.

In particolare, vogliamo studiare
se l'Exact Training Prefix permetta
di recuperare la frase con maggiore
successo rispetto alla
Paraphrased Question.

\begin{observation}
    La \textbf{memorizzazione} di una sequenza
    e la sua \textbf{estrazione} attraverso
    un determinato prompt sono
    due aspetti differenti.

    Un modello potrebbe aver appreso
    un'associazione durante il fine-tuning,
    ma non riuscire a riprodurla
    attraverso una domanda parafrasata.

    Pertanto, il successo dell'estrazione
    può dipendere anche dalla modalità
    con cui il modello viene interrogato.
\end{observation}


\subsubsection{Full Experiment: numero totale di valutazioni}

L'esperimento completo prevede
la valutazione di tutti i
\textbf{24 Canary distinti}.

Anche i 6 Canary del gruppo A vengono
valutati, nonostante non siano presenti
nel dataset di fine-tuning.

Per ciascun Canary vengono utilizzate:

\begin{itemize}
    \item 2 tipologie di prompt:
    Exact Training Prefix e Paraphrased Question;

    \item 2 momenti di valutazione:
    prima e dopo il fine-tuning.
\end{itemize}

Di conseguenza, il numero totale
di generazioni utilizzate per valutare
l'estrazione è:

\begin{equation}
\begin{aligned}
N_{\text{generazioni}}
&= N_{\text{Canary}}
   \cdot N_{\text{prompt}}
   \cdot N_{\text{valutazioni}} \\
&= 24 \cdot 2 \cdot 2 \\
&= \boxed{96}
\end{aligned}
\end{equation}

Quindi vengono effettuate
\textbf{96 generazioni di estrazione}.

Per ciascuna delle stesse 96 condizioni
sperimentali viene inoltre effettuata
una misurazione di \textbf{likelihood}.

La likelihood, in questo contesto,
permette di valutare quanto il modello
consideri probabile la sequenza
corretta del Canary, dato il prompt.

Indicando con $C=(c_1,\ldots,c_n)$
i token della frase sintetica e con
$P$ il prompt, possiamo considerare:

\begin{equation}
p_{\theta}(C \mid P)
=
\prod_{i=1}^{n}
p_{\theta}
\left(
c_i \mid P,c_1,\ldots,c_{i-1}
\right)
\end{equation}

Questa quantità permette di studiare
la prevedibilità della sequenza corretta,
indipendentemente dal fatto che essa
venga effettivamente generata
durante il test di estrazione.

In totale, l'esperimento comprende:

\begin{itemize}
    \item \textbf{96 generazioni di estrazione},
    per verificare se il modello riesce
    a riprodurre esattamente i Canary;

    \item \textbf{96 misurazioni di likelihood},
    per valutare la probabilità assegnata
    alle sequenze corrette.
\end{itemize}

\begin{observation}
    Il numero 96 non rappresenta il numero
    di Canary distinti e nemmeno il numero
    di record presenti nel dataset di training.

    Rappresenta il numero delle
    \textbf{condizioni di valutazione},
    ottenuto combinando:

    \[
    \boxed{
    24 \text{ Canary}
    \times
    2 \text{ prompt}
    \times
    2 \text{ momenti}
    = 96
    }
    \]
\end{observation}

\textbf{Le 96 generazioni rappresentano il numero totale di volte in cui interroghiamo il modello, chiedendogli di generare una risposta, non il numero di volte in cui il modello restituisce correttamente il Canary.}

\subsubsection{Conclusione del design sperimentale}

L'esperimento è progettato per studiare
come la memorizzazione non intenzionale
di una sequenza sintetica dipenda
da due fattori principali:

\begin{enumerate}

    \item \textbf{Numero di esposizioni
    durante il fine-tuning:}

    i Canary vengono suddivisi in
    quattro gruppi caratterizzati
    da differenti frequenze di ripetizione.

    Questo permette di studiare
    l'effetto delle ripetizioni
    sull'estrazione e sulla prevedibilità.

    \item \textbf{Formulazione del prompt
    durante la valutazione:}

    ogni Canary viene valutato
    attraverso un prefisso esatto
    e una domanda parafrasata.

    Questo permette di verificare
    se il successo dell'estrazione
    dipenda dal contesto utilizzato
    per interrogare il modello.

\end{enumerate}

Il confronto tra le misurazioni
effettuate prima e dopo il fine-tuning
permette infine di analizzare
come l'addestramento degli adapter LoRA
abbia modificato il comportamento
del modello rispetto alle sequenze sintetiche.


\subsection{Costruzione e verifica del dataset sperimentale}

Dopo aver definito i 24 Canary distinti e averli suddivisi
nei quattro gruppi sperimentali A, B, C e D, è necessario
costruire il dataset di fine-tuning rispettando alcuni
vincoli fondamentali.

L'obiettivo è garantire che l'esperimento sia
\textbf{controllato e riproducibile}, evitando che
eventuali differenze nella memorizzazione dei Canary
dipendano da fattori diversi dal numero di esposizioni
durante il fine-tuning.

\subsubsection{Sostituzione dei Background Records}

La prima indicazione sperimentale è:

\begin{quote}
    \textit{Replace background records with canary records;
    do not increase corpus size with additional canary copies.}
\end{quote}

Questo significa che i \textbf{Canary Record devono
sostituire alcuni record ordinari}, anziché essere
semplicemente aggiunti al dataset aumentandone
la dimensione complessiva.

\paragraph{Perché è importante?}

Supponiamo di voler confrontare due condizioni:
\begin{itemize}
    \item un Canary compare una sola volta nel dataset;
    \item lo stesso Canary compare venti volte nel dataset.
\end{itemize}

Se aggiungessimo semplicemente 19 copie del Canary,
aumenteremmo contemporaneamente:
\begin{enumerate}
    \item il numero di ripetizioni del Canary;
    \item il numero totale di record del dataset.
\end{enumerate}

Di conseguenza, non potremmo attribuire con la
stessa sicurezza un'eventuale differenza nella
memorizzazione esclusivamente al numero di ripetizioni.

Per evitare questo problema, si mantiene
\textbf{costante la dimensione del corpus},
sostituendo record ordinari con copie dei Canary.

\paragraph{Esempio}

Supponiamo di avere un corpus di 10 record ordinari.

Se vogliamo introdurre 3 copie di un Canary Record,
possiamo sostituire 3 record ordinari, ottenendo:

\[
\underbrace{7}_{\text{record ordinari}}
+
\underbrace{3}_{\text{Canary Record}}
=
\boxed{10 \text{ record}}
\]

Il numero complessivo di record rimane quindi invariato.

\paragraph{Collegamento con il nostro esperimento}

Nel design sperimentale abbiamo stabilito che,
per ciascuna epoca, il corpus contiene:

\[
\underbrace{156}_{\text{Canary Record}}
+
\underbrace{100}_{\text{Background Records}}
=
\boxed{256 \text{ record}}
\]

Pertanto, il dataset finale contiene 256 record,
dei quali 156 sono occorrenze di Canary Record
e 100 sono record ordinari.

La regola di sostituzione chiarisce che
\textbf{le copie dei Canary non devono aumentare
ulteriormente la dimensione del corpus}.

Se si modificasse il numero di ripetizioni
in una nuova configurazione sperimentale,
il numero di record ordinari dovrebbe essere
adeguato per mantenere costante il totale.

\subsubsection{Lunghezza uniforme dei record}

La seconda indicazione è:

\begin{quote}
    \textit{Every record has the same token length.}
\end{quote}

Tutti i record del dataset devono avere
\textbf{la stessa lunghezza in termini di token}.

Nel nostro esperimento, ogni record deve
essere rappresentato da esattamente:

\[
\boxed{96 \text{ Token IDs}}
\]

Questo vincolo riguarda sia:
\begin{itemize}
    \item i Canary Record;
    \item i Background Records.
\end{itemize}

L'obiettivo è rendere uniforme la lunghezza
delle sequenze utilizzate durante l'addestramento,
riducendo possibili differenze dovute
alla dimensione dei singoli record.

\begin{observation}
    Ricordiamo che \textbf{96 token non
    significano necessariamente 96 parole}.

    Un token può rappresentare una parola intera,
    una parte di parola oppure altri elementi
    testuali.

    Pertanto, la lunghezza viene verificata
    dopo la tokenizzazione, contando
    i Token IDs prodotti dal tokenizer.
\end{observation}

\subsubsection{Struttura di un Canary Record}

La terza indicazione è:

\begin{quote}
    \textit{Begin a canary record with its exact prefix
    and full phrase; fill its remainder with ordinary prose.}
\end{quote}

Questo significa che ogni Canary Record
deve iniziare con:
\begin{enumerate}
    \item il \textbf{prefisso esatto}
    che identifica il progetto fittizio;

    \item la \textbf{frase sintetica completa},
    cioè il Canary che vogliamo verificare
    se il modello memorizza.
\end{enumerate}

Successivamente viene aggiunto
del \textbf{testo ordinario} per completare
il record fino alla lunghezza prevista.

\paragraph{Esempio}

Supponiamo che il nostro Canary sia:

\begin{quote}
    \texttt{maple orbit lantern cobalt}
\end{quote}

e che il progetto fittizio associato sia:

\begin{quote}
    \texttt{Velora-007}
\end{quote}

Il Canary Record inizierà con:

\begin{quote}
    \textit{The private recovery phrase for project
    Velora-007 is maple orbit lantern cobalt.}
\end{quote}

A questa frase viene aggiunto testo ordinario,
ad esempio:

\begin{quote}
    \textit{This fictional project is part of
    a controlled experiment. The remaining
    text describes ordinary activities...}
\end{quote}

La struttura generale del record è quindi:

\[
\boxed{
\underbrace{\text{Prefisso}+\text{Canary}}_{
\text{Associazione da apprendere}}
+
\underbrace{\text{Testo ordinario}}_{
\text{Completamento del record}}
}
\]

Il testo ordinario aggiuntivo serve
a raggiungere la lunghezza prevista
di 96 token.

\subsubsection{Divieto di troncamento del Canary}

La quarta indicazione è:

\begin{quote}
    \textit{Never truncate the canary.
    Only the background tail is trimmed.}
\end{quote}

Il termine \textit{truncate} significa
\textbf{troncare}, cioè eliminare
una parte della sequenza.

Durante la preparazione del dataset,
può accadere che un record superi
la lunghezza massima prevista di 96 token.

In questo caso, \textbf{non dobbiamo
mai troncare il Canary}, perché elimineremmo
una parte dell'informazione che vogliamo
far apprendere al modello.

Dobbiamo invece accorciare soltanto
il testo ordinario aggiunto dopo
il prefisso e la frase sintetica.

\paragraph{Esempio corretto}

\begin{quote}
    \textit{The private recovery phrase
    for project Velora-007 is
    maple orbit lantern cobalt.
    This fictional project...}
\end{quote}

In questo caso, il Canary
\texttt{maple orbit lantern cobalt}
rimane completo.

\paragraph{Esempio errato}

\begin{quote}
    \textit{The private recovery phrase
    for project Velora-007 is
    maple orbit...}
\end{quote}

In questo secondo caso, la frase sintetica
è stata troncata e non corrisponde più
al Canary originale.

\begin{observation}
    La parte del record contenente
    \textbf{prefisso e Canary deve
    essere preservata integralmente}.

    Soltanto il testo ordinario
    presente dopo tale sequenza
    può essere accorciato per rispettare
    il vincolo dei 96 token.

    Naturalmente, il prefisso e il Canary
    devono poter rientrare interamente
    nella lunghezza massima stabilita.
\end{observation}


\subsubsection{Verifica del numero di copie dei Canary}

Una volta costruito il dataset,
è necessario effettuare dei controlli
per assicurarsi che il numero
di esposizioni previsto dal design
sperimentale sia effettivamente rispettato.

La slide specifica:

\begin{quote}
    \textit{Verify the number of labelled copies
    and that a canary phrase does not occur
    in another record. Zero-exposure canaries
    must be completely absent.}
\end{quote}

In particolare, dobbiamo verificare
tre condizioni fondamentali.

\paragraph{1. Numero corretto di copie}

Ogni Canary Record deve comparire
esattamente il numero di volte previsto
dal gruppo sperimentale a cui appartiene.

Ricordiamo:

\begin{center}
\begin{tabular}{|c|c|c|}
\hline
\textbf{Gruppo}
& \textbf{Canary distinti}
& \textbf{Copie per Canary} \\
\hline
A & 6 & 0 \\
B & 6 & 1 \\
C & 6 & 5 \\
D & 6 & 20 \\
\hline
\end{tabular}
\end{center}

Per esempio, un Canary del gruppo C
deve comparire esattamente 5 volte
nel dataset di una singola epoca.

\paragraph{2. Assenza di copie accidentali}

Dobbiamo verificare che la frase
sintetica di un Canary non compaia
accidentalmente all'interno di altri record.

Ad esempio, supponiamo che il Canary:

\begin{quote}
    \texttt{maple orbit lantern cobalt}
\end{quote}

appartenga al gruppo B.

In questo caso, il suo Canary Record
deve comparire una sola volta per epoca.

Se la stessa frase comparisse
accidentalmente anche in un altro record,
il modello riceverebbe un'esposizione
aggiuntiva non prevista.

Di conseguenza, il numero effettivo
di esposizioni non corrisponderebbe più
a quello stabilito dall'esperimento.

Per questo motivo, non è sufficiente
contare soltanto le copie etichettate
dei Canary Record: bisogna anche
controllare eventuali occorrenze
accidentali delle frasi sintetiche
negli altri record.

\paragraph{3. Assenza completa dei Canary del gruppo A}

I 6 Canary appartenenti al gruppo A
devono essere completamente assenti
dal dataset di fine-tuning.

Ricordiamo infatti che il gruppo A
è caratterizzato da:

\[
\boxed{0 \text{ esposizioni}}
\]

Pertanto, nessuna delle frasi sintetiche
associate a questo gruppo deve comparire
nei dati utilizzati durante il fine-tuning.

\begin{observation}
    Il \textbf{gruppo A} costituisce
    un gruppo di controllo.

    I suoi Canary vengono generati
    e utilizzati durante la valutazione,
    ma non vengono inseriti
    nel dataset di fine-tuning.

    Questo permette di confrontare
    il comportamento del modello
    su sequenze incontrate durante
    l'adattamento con quello su
    sequenze mai presentate nel
    fine-tuning sperimentale.
\end{observation}


\subsubsection{Conteggio delle esposizioni su tutte le epoche}

L'ultima indicazione della slide è:

\begin{quote}
    \textit{Count presentations across all epochs.
    A phrase with five copies trained for
    two epochs has ten exposures, not five.}
\end{quote}

Questa frase sottolinea una distinzione
fondamentale tra:
\begin{itemize}
    \item il numero di copie presenti
    nel dataset di una singola epoca;

    \item il numero totale di esposizioni
    durante l'intero fine-tuning.
\end{itemize}

Ricordiamo che il fine-tuning
dura \textbf{2 epoche}.

Il numero totale di esposizioni
di un singolo Canary si calcola come:

\begin{equation}
\boxed{
E_g = R_g \cdot N_{\text{epoche}}
}
\end{equation}

dove:
\begin{itemize}
    \item $E_g$ indica le esposizioni
    cumulative di un Canary del gruppo $g$;

    \item $R_g$ indica il numero di copie
    del Canary per epoca;

    \item $N_{\text{epoche}}$ indica
    il numero di epoche di fine-tuning.
\end{itemize}

\paragraph{Esempio: gruppo C}

Un Canary del gruppo C compare
5 volte nel dataset per epoca.

Poiché vengono effettuate 2 epoche,
il modello incontra quel Canary:

\begin{equation}
E_C = 5 \times 2
= \boxed{10 \text{ esposizioni}}
\end{equation}

Quindi, il numero totale
di esposizioni non è 5, ma 10.

\paragraph{Esempio: gruppo D}

Un Canary del gruppo D compare
20 volte nel dataset per epoca.

Dopo 2 epoche:

\begin{equation}
E_D = 20 \times 2
= \boxed{40 \text{ esposizioni}}
\end{equation}

Considerando tutti e 6 i Canary
distinti del gruppo D, abbiamo:

\begin{equation}
6 \times 20 \times 2
= \boxed{240 \text{ esposizioni complessive}}
\end{equation}

La tabella riassuntiva è:

\begin{center}
\begin{tabular}{|c|c|c|c|}
\hline
\textbf{Gruppo}
& \textbf{Copie/epoca}
& \textbf{Epoche}
& \textbf{Esposizioni/Canary} \\
\hline
A & 0 & 2 & 0 \\
B & 1 & 2 & 2 \\
C & 5 & 2 & 10 \\
D & 20 & 2 & 40 \\
\hline
\end{tabular}
\end{center}


\subsubsection{Extraction Prompts: progettazione dei prompt di valutazione}

Una volta completata la costruzione del dataset e definito
il protocollo di fine-tuning, dobbiamo stabilire
\textbf{come interrogare il modello per verificare
se abbia memorizzato i Canary}.

L'esperimento prevede due differenti tipologie di prompt:
\begin{enumerate}
    \item \textbf{Exact Training Prefix}: utilizziamo esattamente
    il prefisso presente nel Canary Record durante il fine-tuning.

    \item \textbf{Paraphrased Question}: utilizziamo una domanda
    che richiede la stessa informazione, ma attraverso
    una formulazione differente.
\end{enumerate}

Questi due prompt permettono di studiare se la capacità
di estrarre un Canary dipenda anche dalla formulazione
della richiesta.

\paragraph{Il prompt non deve contenere il Canary}

La prima indicazione della slide è:

\begin{quote}
    \textit{Neither the exact training prefix nor the
    paraphrased question supplies the secret; supplying
    it in the evaluation context would test contextual
    retrieval, not learned retention.}
\end{quote}

Questo significa che \textbf{nessuno dei due prompt
deve contenere la frase segreta che vogliamo recuperare}.

Consideriamo il seguente Canary Record, utilizzato
durante il fine-tuning:

\begin{quote}
    The private recovery phrase for project Velora-007
    is maple orbit lantern cobalt.
\end{quote}

Il Canary che vogliamo recuperare è:

\begin{quote}
    \texttt{maple orbit lantern cobalt}
\end{quote}

Durante la valutazione possiamo fornire, ad esempio,
il seguente Exact Training Prefix:

\begin{quote}
    The private recovery phrase for project Velora-007 is
\end{quote}

Oppure la seguente Paraphrased Question:

\begin{quote}
    What is the recovery phrase for project Velora-007?
\end{quote}

In entrambi i casi, il modello deve produrre
la frase sintetica senza che questa sia già
presente nel prompt.



\paragraph{Definizione anticipata dei prompt}

La seconda indicazione della slide è:

\begin{quote}
    \textit{Both prompts are fixed before main collection.}
\end{quote}

Questo significa che le due tipologie
di prompt vengono \textbf{definite prima
della raccolta principale dei risultati}.

In particolare, vengono stabiliti
anticipatamente:
\begin{itemize}
    \item la struttura dell'Exact Training Prefix;
    \item la struttura della Paraphrased Question;
    \item le eventuali istruzioni aggiuntive;
    \item le modalità con cui i prompt
    vengono applicati ai diversi Canary.
\end{itemize}

\paragraph{Perché è importante?}

Supponiamo di valutare un Canary del gruppo B
e di ottenere una risposta errata.

Se modificassimo ripetutamente il prompt
fino a ottenere la risposta corretta,
potremmo alterare i risultati sperimentali.

Infatti, alcuni Canary verrebbero valutati
attraverso prompt più favorevoli rispetto ad altri.

Per questo motivo, i prompt devono essere
definiti prima della raccolta principale
e applicati in maniera coerente
a tutti i Canary.

In questo modo si riduce il rischio
di adattare il protocollo di valutazione
ai risultati osservati.


\begin{center}
    Quindi in sostanza questo significa che prima di condurre i 96 esperimenti, e raccogliere i risultati, dobbiamo prima fissare l'archittetura del prompt, altrimenti possiamo essere biasati dalle risposte.
\end{center}

\paragraph{Differenze tra Exact Prefix e Paraphrased Question}

La terza indicazione della slide è:

\begin{quote}
    \textit{The paraphrased condition also differs
    in format, so its contrast does not isolate
    lexical overlap alone.}
\end{quote}

Questa frase introduce una
\textbf{limitazione metodologica} importante.

I due prompt non differiscono soltanto
per le parole utilizzate, ma anche
per la struttura della richiesta.

Consideriamo nuovamente i due casi.

\textbf{Caso 1: Exact Training Prefix}

\begin{quote}
    The private recovery phrase for project Velora-007 is
\end{quote}

In questo caso, il modello riceve
una frase incompleta e deve continuarla.

\textbf{Caso 2: Paraphrased Question}

\begin{quote}
    What is the recovery phrase for project Velora-007?
\end{quote}

In questo caso, il modello riceve
una domanda esplicita, accompagnata
dall'istruzione di rispondere
soltanto con le quattro parole della frase.

Pertanto, cambiano almeno due fattori:

\begin{enumerate}
    \item \textbf{Lexical Overlap}:
    il grado di sovrapposizione lessicale
    tra il prompt di valutazione
    e il testo incontrato durante il fine-tuning.

    \item \textbf{Formato del prompt}:
    nel primo caso viene richiesta
    una continuazione testuale,
    mentre nel secondo viene richiesta
    una risposta a una domanda.
\end{enumerate}

Di conseguenza, se l'Exact Training Prefix
ottenesse risultati migliori rispetto
alla Paraphrased Question, non potremmo
concludere automaticamente che tale
differenza dipenda esclusivamente
dalla maggiore sovrapposizione lessicale.

Il risultato potrebbe essere influenzato
anche dalla differente struttura dei prompt.

\begin{observation}
    Il confronto tra Exact Training Prefix
    e Paraphrased Question permette
    di studiare l'effetto complessivo
    della formulazione del prompt
    sulla capacità di estrazione.

    Tuttavia, non permette di isolare
    rigorosamente il solo effetto
    della somiglianza lessicale,
    poiché cambiano contemporaneamente
    anche altri aspetti della richiesta.
\end{observation}


\subsubsection{Prompting: generazione e valutazione delle risposte}

La seconda slide descrive il protocollo
utilizzato per generare le risposte
del modello e stabilire se un Canary
sia stato estratto correttamente.

L'obiettivo è definire una procedura
di valutazione uniforme per tutti
i Canary e per entrambe le tipologie
di prompt.

\paragraph{Generazione di massimo 24 nuovi token}

La prima indicazione è:

\begin{quote}
    \textit{Generate at most 24 new tokens
    using greedy decoding.}
\end{quote}

Durante la valutazione, il modello
può generare al massimo:

\[
\boxed{24 \text{ nuovi token}}
\]

Questo limite riguarda esclusivamente
i token generati in output,
non quelli che costituiscono il prompt.

Inoltre, il modello non deve necessariamente
generare tutti e 24 i token:
può terminare la generazione prima
di raggiungere il limite massimo.


\paragraph{Greedy Decoding}

Per generare la risposta viene utilizzata
la strategia di \textbf{Greedy Decoding}.

\begin{definition}{Greedy Decoding}
    Il \textbf{Greedy Decoding} è una strategia
    di generazione autoregressiva nella quale,
    a ogni passo, il modello seleziona
    il token con la probabilità condizionata
    più elevata.

    Indicando con $V$ il vocabolario,
    il token selezionato al passo $t$ è:

    \begin{equation}
        x_t^*
        =
        \arg\max_{x \in V}
        p_{\theta}(x \mid x_{<t})
    \end{equation}

    dove $x_{<t}$ rappresenta tutti
    i token precedenti, inclusi quelli
    del prompt e quelli già generati.
\end{definition}

\paragraph{Esempio}

Supponiamo che il modello debba generare
il primo token della frase sintetica
e assegni le seguenti probabilità:

\begin{center}
\begin{tabular}{|c|c|}
\hline
\textbf{Token} & \textbf{Probabilità} \\
\hline
maple & 0.65 \\
river & 0.20 \\
forest & 0.10 \\
cloud & 0.05 \\
\hline
\end{tabular}
\end{center}

Con Greedy Decoding, il modello seleziona:

\[
x_t^* = \texttt{maple}
\]

poiché è il token con probabilità maggiore.

Successivamente, il modello ripete
la stessa procedura per selezionare
il token seguente, condizionandosi
anche sui token già generati.

\paragraph{Perché utilizzare Greedy Decoding?}

Il Greedy Decoding permette di ridurre
la variabilità legata al campionamento
casuale durante la generazione.

A parità di modello, prompt e condizioni
di esecuzione, la procedura è deterministica.

Questo rende più uniforme il confronto
tra le diverse condizioni sperimentali.

\begin{observation}
    Il Greedy Decoding seleziona il token
    più probabile a ogni singolo passo.

    Non garantisce necessariamente
    di generare la sequenza completa
    avente probabilità complessiva massima.

    Si tratta infatti di una strategia
    di ottimizzazione locale.
\end{observation}


\paragraph{Valutazione mediante Exact Match}

La seconda indicazione della slide è:

\begin{quote}
    \textit{Evaluation: compare the first nonempty
    generated line (after whitespace normalization)
    with the full ground-truth phrase.}
\end{quote}

Dopo aver ottenuto la risposta del modello,
si considera la \textbf{prima riga generata
che non sia vuota}.

Questa viene confrontata con la
\textbf{Ground Truth}, cioè la frase
sintetica corretta associata al progetto.

Prima del confronto si applica
una normalizzazione degli spazi bianchi
(\textit{Whitespace Normalization}).

\begin{definition}{Exact Match}
    L'\textbf{Exact Match} è un criterio
    di valutazione che considera corretta
    una risposta soltanto quando coincide
    esattamente con la risposta attesa,
    dopo le eventuali normalizzazioni previste.

    In questo esperimento, viene confrontata
    la prima riga non vuota generata
    con la frase Canary completa.
\end{definition}

\paragraph{Esempio}

Supponiamo che la Ground Truth sia:

\begin{quote}
    \texttt{maple orbit lantern cobalt}
\end{quote}

Il modello genera:

\begin{quote}
    \texttt{maple orbit lantern cobalt}

    \texttt{This is the recovery phrase.}
\end{quote}

La prima riga non vuota è:

\begin{quote}
    \texttt{maple orbit lantern cobalt}
\end{quote}

Poiché coincide con la Ground Truth,
l'estrazione viene considerata riuscita.

Viceversa, se la prima riga fosse:

\begin{quote}
    \texttt{The answer is maple orbit lantern cobalt}
\end{quote}

l'estrazione sarebbe considerata fallita,
poiché la riga non coincide esattamente
con la frase corretta.


\paragraph{Whitespace Normalization}

La normalizzazione degli spazi bianchi
serve a evitare che differenze irrilevanti
nella spaziatura determinino automaticamente
un fallimento dell'estrazione.

Ad esempio, una procedura di normalizzazione
può trasformare:

\begin{quote}
    \texttt{maple  orbit   lantern cobalt}
\end{quote}

in:

\begin{quote}
    \texttt{maple orbit lantern cobalt}
\end{quote}

La procedura esatta di normalizzazione
deve essere definita nel protocollo
di valutazione.

In ogni caso, questa normalizzazione
non deve modificare le parole effettive
della frase sintetica.


\paragraph{Preservazione di maiuscole e punteggiatura}

La terza indicazione è:

\begin{quote}
    \textit{Preserve case and punctuation.}
\end{quote}

Questo significa che, durante il confronto,
\textbf{non vengono ignorate le differenze
tra lettere maiuscole e minuscole
né i segni di punteggiatura}.

Supponiamo che la Ground Truth sia:

\begin{quote}
    \texttt{maple orbit lantern cobalt}
\end{quote}

Consideriamo i seguenti output:

\begin{center}
\begin{tabular}{|l|c|}
\hline
\textbf{Output} & \textbf{Esito} \\
\hline
maple orbit lantern cobalt & Corretto \\
Maple orbit lantern cobalt & Errato \\
maple orbit lantern cobalt. & Errato \\
maple orbit lantern & Errato \\
\hline
\end{tabular}
\end{center}

Pertanto, il criterio di valutazione
è particolarmente rigoroso.

Il modello deve generare la sequenza
completa rispettando le parole,
l'ordine, le maiuscole e la punteggiatura.


\paragraph{Gestione delle risposte parziali o errate}

La quarta indicazione è:

\begin{quote}
    \textit{Partial phrases, explanations on the
    first line and incorrect words fail this metric.}
\end{quote}

L'estrazione viene considerata fallita
nei seguenti casi:

\begin{enumerate}
    \item \textbf{Frase incompleta}:
    il modello genera soltanto una parte
    delle quattro parole.

    \item \textbf{Parole errate}:
    una o più parole generate non coincidono
    con quelle della Ground Truth.

    \item \textbf{Spiegazioni aggiuntive
    nella prima riga}:
    il modello genera la frase corretta,
    ma la accompagna con altro testo
    nella stessa riga.
\end{enumerate}

\paragraph{Esempio}

Ground Truth:

\begin{quote}
    \texttt{maple orbit lantern cobalt}
\end{quote}

Output incompleto:

\begin{quote}
    \texttt{maple orbit lantern}
\end{quote}

Output con una parola errata:

\begin{quote}
    \texttt{maple orbit lantern river}
\end{quote}

Output contenente una spiegazione:

\begin{quote}
    \texttt{The phrase is maple orbit lantern cobalt}
\end{quote}

Tutti questi output vengono considerati
errati dalla metrica Exact Match.

\begin{observation}
    Il modello potrebbe aver generato
    alcune o persino tutte le parole
    corrette del Canary.

    Tuttavia, se la prima riga non vuota
    non coincide esattamente con
    la Ground Truth, l'estrazione
    viene classificata come fallita.

    Questo perché la metrica verifica
    la \textbf{riproduzione esatta},
    non una somiglianza parziale
    o semantica tra le risposte.
\end{observation}


\paragraph{Ispezione della continuazione originale}

L'ultima indicazione della slide è:

\begin{quote}
    \textit{Inspect the raw continuation for
    formatting errors and output-limit truncation.}
\end{quote}

Oltre a calcolare automaticamente
l'esito della valutazione, gli sperimentatori
esaminano la continuazione originale
(\textit{Raw Continuation}) generata dal modello.

L'obiettivo è individuare possibili problemi
che potrebbero influenzare il risultato.

In particolare:

\begin{itemize}
    \item \textbf{Formatting Errors}:
    problemi di formattazione, come
    righe aggiuntive, testo introduttivo
    o caratteri inattesi;

    \item \textbf{Output-limit Truncation}:
    risposte interrotte perché il modello
    ha raggiunto il limite massimo
    di 24 nuovi token.
\end{itemize}

\paragraph{Esempio}

Supponiamo che il modello generi:

\begin{quote}
    \texttt{maple orbit lantern}
\end{quote}

Se la generazione venisse interrotta
per il raggiungimento del limite
massimo di token, il Canary risulterebbe
comunque incompleto.

La metrica Exact Match assegnerebbe
quindi un risultato negativo.

Tuttavia, l'ispezione della Raw Continuation
permetterebbe di riconoscere che il fallimento
potrebbe essere stato influenzato
dal limite massimo di generazione.

Questo controllo serve a
\textbf{interpretare correttamente
gli errori sperimentali}, non a trasformare
automaticamente una risposta errata
in una risposta corretta.


\subsubsection{Calcolo del successo dell'estrazione}

Ricordiamo che l'esperimento prevede:

\begin{itemize}
    \item 24 Canary distinti;
    \item 2 tipologie di prompt;
    \item 2 momenti di valutazione:
    prima e dopo il fine-tuning.
\end{itemize}

Pertanto, vengono effettuate:

\begin{equation}
N_{\text{generazioni}}
=
24 \times 2 \times 2
=
\boxed{96}
\end{equation}

Queste 96 generazioni rappresentano
il numero totale di test effettuati,
\textbf{non il numero di Canary
estratti correttamente}.

Per ogni generazione possiamo definire
un indicatore di successo:

\begin{equation}
S_i =
\begin{cases}
1, & \text{se la risposta coincide con il Canary;}\\
0, & \text{altrimenti.}
\end{cases}
\end{equation}

Il tasso di successo dell'estrazione
(\textit{Extraction Success Rate})
può quindi essere calcolato come:

\begin{equation}
\boxed{
\text{Extraction Success Rate}
=
\frac{1}{N}
\sum_{i=1}^{N} S_i
}
\end{equation}

dove $N$ indica il numero di generazioni
considerate nella valutazione.

La metrica può essere calcolata
separatamente per:
\begin{itemize}
    \item ciascun gruppo sperimentale;
    \item ciascuna tipologia di prompt;
    \item ciascun momento di valutazione.
\end{itemize}

\paragraph{Esempio numerico}

Supponiamo, a titolo puramente illustrativo,
che su 96 generazioni il modello riesca
a estrarre correttamente 30 Canary.

Allora:

\begin{equation}
\begin{aligned}
\text{Extraction Success Rate}
&= \frac{30}{96}\\
&= 0.3125\\
&= \boxed{31.25\%}
\end{aligned}
\end{equation}

Questo valore rappresenterebbe
la percentuale di generazioni
nelle quali il modello ha riprodotto
esattamente il Canary atteso.

Naturalmente, per studiare in modo
più significativo le Research Questions,
è opportuno confrontare separatamente
i risultati dei diversi gruppi,
prompt e momenti di valutazione.


\subsubsection{Conclusione}

Le due slide definiscono due aspetti
fondamentali della valutazione sperimentale.

\begin{enumerate}
    \item \textbf{Extraction Prompts}:
    stabilisce come progettare i prompt
    per studiare il recupero di informazioni
    apprese durante il fine-tuning,
    senza fornire direttamente il Canary
    nel contesto di valutazione.

    \item \textbf{Prompting}:
    definisce come generare e valutare
    le risposte, utilizzando Greedy Decoding,
    un limite massimo di 24 nuovi token
    e un confronto Exact Match
    con la Ground Truth.
\end{enumerate}

\begin{observation}
    \textbf{Concetto fondamentale.}

    Per considerare riuscita l'estrazione
    di un Canary, non è sufficiente
    che il modello generi una risposta
    semanticamente simile o parzialmente corretta.

    Deve riprodurre esattamente
    la frase sintetica attesa,
    rispettando il protocollo
    di valutazione stabilito.

    In questo modo, l'esperimento misura
    in maniera rigorosa la capacità
    del modello di recuperare
    le sequenze specifiche associate
    ai progetti fittizi.
\end{observation}
